% -----------------------------------------------------------
\section{Merit}
% -----------------------------------------------------------
	\label{MERIT}
	
% % ----------------------
% \subsection{see also}
% % ----------------------

% ----------------------
\subsection{1828 American Dictionary of the English Language} % *1828
% ----------------------
	\label{MERIT-1828}
	
	% -----
	\subsubsection{Noun}
	% -----

		\begin{compactitem}
			\item[1] Desert; goodness or excellence which entitles one to honor or regard; worth; any performance or worth which claims regard or compensation; applied to morals, to excellence in writing, or to valuable services of any kind. Thus we speak of the inability of men to obtain salvation by their own merits. We speak of the merits of an author; the merits of a soldier, etc.
			\item[2] Value; excellence; applied to things; as the merits of an essay or poem; the merits of a pointing; the merits of a heroic achievement.
			\item[3] Reward deserved; that which is earned or merited.
		\end{compactitem}
		
	% -----
	\subsubsection{Verb}
	% -----
	
		\begin{compactitem}
			\item To deserve; to earn by active service, or by any valuable performance; to have a right to claim reward in money, regard, honor or happiness. Watts, by his writings, merited the gratitude of the whole christian world. The faithful laborer merits his wages.
			\item[1] To deserve; to have a just title to. Fidelity merits and usually obtains confidence.
			\item[2] To deserve, in an ill sense; to have a just title to. Every violation of law merits punishment. Every sin merits God's displeasure.
		\end{compactitem} 

% ----------------------
\subsection{Oxford English Dictionary} % *OED
% ----------------------
	\label{MERIT-OED}
	
	% -----
	\subsubsection{Noun}
	% -----

		\begin{compactitem}
			\item[I] As an abstract quality.
			\item[I.1.a] \textit{Theology.} The quality (in actions or persons) of being entitled to reward from God.
			\item[I.1.b] \textit{gen.} The quality of deserving well, or of being entitled to reward or gratitude.
			\item[I.1.c] Claim or title to commendation or esteem; excellence, worth.
			\item[I.1.d] The condition of being valued or honoured; esteem. \textit{Obsolete.}
			\item[I.2] The condition or fact of deserving reward or punishment. \textit{Obsolete.}
			\item[I.3] Claim to gratitude for being the agent of some favourable state of affairs; the honour or credit of bringing about (something). With \textit{of}. Now \textit{rare}.
			\item[II] something which confers the quality of merit, and other more concrete senses.
			\item[II.4.a] \textit{Theology.} In \textit{plural.} Good works viewed as entitling a person to reward from God; (also) the righteousness and sacrifice of Christ as the ground on which God grants forgiveness to sinners. Also (occasionally) in \textit{singular.}
				\begin{compactitem}
					\item[1662] ``That we fail not finally to attain thy heavenly promises, through the merits of Jesus Christ our Lord'' (\textit{Book of Common Prayer}).
				\end{compactitem}
			\item[II.4.b] In \textit{plural.} That which entitles a person to reward or gratitude. In early use also more generally: deserts. Now merged in sense II.4c. \textit{Obsolete.}
			\item[II.4.c] A commendable quality, an excellence, a good point. Now frequently in \textbf{to have the merit that} (also \textbf{of}, etc.).
			\item[II.4.d] \textit{Buddhism} and \textit{Jainism}. The good actions which a person performs in one state of existence, which help determine his or her successive entry into a better state; the effect on a person's karma of unselfish actions; (also) an unselfish action. Esp. in \textbf{to acquire merit}.
			\item[II.5] That which is deserved or has been earned, whether good or evil; due reward or punishment. \textit{Obsolete.}
			\item[II.6] \textbf{the merits} (rarely, \textbf{the merit}) (of a case, question, etc.): (\textit{Law}) the intrinsic rights and wrongs of a case, in contradistinction to extraneous or technical points, esp. of procedure; (\textit{gen.}) the intrinsic rights and wrongs or excellences and defects of something.
		\end{compactitem}
		
	% -----
	\subsubsection{Verb}
	% -----
	
		\begin{compactitem}
			\item[1] \textit{transitive.} To reward, to recompense. \textit{Obsolete.} \textit{rare.}
			\item[2] To be or become entitled to or worthy of (reward, punishment, good or evil fortune or estimation, etc.)
			\item[2.a] \textit{transitive.} With noun, noun phrase, or gerund as object.
			\item[2.b] \textit{transitive.} With infinitive or (rarely) that-clause as object. In early use (occasionally): to obtain as one's deserts (\textit{obsolete}). Now \textit{archaic}.
			\item[2.c] \textit{transitive.} In clauses with \textit{as} or \textit{than}.
			\item[3] \textit{intransitive.} To acquire merit; to become entitled to reward, gratitude, or commendation. \textit{Obsolete.}
			\item[4] \textit{transitive.} To earn by meritorious action; (\textit{Theology}) to become entitled to (reward) from God; (of Christ) to obtain by vicarious sacrifice (spiritual blessings) for humankind (cf. \textbf{merit} n. II.4a). \textit{Obsolete.}
			\item[5] \textit{intransitive.} To be deserving of good or evil; = \textbf{deserve} v. 3a. Frequently in to merit well (of a person). \textit{Obsolete} (\textit{archaic} in later use).
		\end{compactitem}
	
% % ----------------------
% \subsection{Buck's Theological Dictionary (1833)} % *BTD
% % ----------------------
	% \label{MERIT-BTD}

	% \begin{compactdesc}
		% \item[PAGE\#]
	% \end{compactdesc}
	
% % ----------------------
% \subsection{Restoration Lexicon} % *RL *REsTORATION LEXICON
% % ----------------------
	% \label{MERIT-OED}

	% \begin{compactitem}
		% \item[]
	% \end{compactitem}

% ----------------------
\subsection{scriptural Usage} % *sCRIPTUREs
% ----------------------
	\label{MERIT-sCRIPTUREs}
	
	\begin{tabular}{ll}
		Total occurrences of ``merit'' (noun) in the scriptures & 6 \\
		Total occurrences of ``merit'' (verb) in the scriptures & 2 \\
		Total occurrences in the scriptures & 8 \\
	\end{tabular}

	% -----
	\subsubsection{Old Testament} % *OT
	% -----
		\label{MERIT-OT}
	
		% --
		\paragraph{KJV}
		% --
			\label{MERIT-OT-KJV}
	
			\begin{tabular}{ll}
				Total occurrences in the OT & 0
			\end{tabular}
			
			% \begin{compactdesc}
				% \item[]
			% \end{compactdesc}
			
		% % --
		% \paragraph{Hebrew Bible}
		% % --
			% \label{MERIT-HEBREW}
		
			% %
			% \subparagraph{\texorpdfstring{\he{sample word}}{}} % paste actual hebrew text here
			% %
				% \label{PAGA} % whatever is in step bible without periods
			
				% \begin{compactdesc}
					% \item[HALOT , PDF ] \textbf{} % Use the 1994 PDF
						% \begin{compactitem}
							% \item[1]
						% \end{compactitem}
					% \item[BDB , PDF ] \textbf{}
						% \begin{compactitem}
							% \item[1]
						% \end{compactitem}
					% \item[DCH PDF ] \textbf{} % don't include the actual page number, they repeat from volume to volume
						% \begin{compactitem}
							% \item[1]
						% \end{compactitem}
				% \end{compactdesc}
				
				% \begin{compactdesc}
					% \item[some OT uses] \textbf{}
						% \begin{compactdesc}
							% \item[]
						% \end{compactdesc}
				% \end{compactdesc}
			
		% % --
		% \paragraph{septuagint}
		% % --
			% \label{MERIT-LXX}
		
			% %
			% \subparagraph{\texorpdfstring{\gr{sample word}}{}}
			% %
		
	% -----
	\subsubsection{New Testament} % *NT
	% -----
		\label{MERIT-NT}
	
		% --
		\paragraph{KJV}
		% --
			\label{MERIT-NT-KJV}		
	
			\begin{tabular}{ll}
				Total occurrences in the NT & 0
			\end{tabular}
			
			% \begin{compactdesc}
				% \item[]
			% \end{compactdesc}
			
		% % --
		% \paragraph{Greek New Testament} % *GREEK
		% % --
			% \label{MERIT-GREEK}
		
			% %
			% \subparagraph{\texorpdfstring{\gr{sample word}}{}}
			% %
				% \label{KATALLAGE}
				
				% \begin{tabular}{ll}
					% Total occurrences in the NT & 
				% \end{tabular}
				
				% \begin{compactdesc}
					% \item[some NT uses] \textbf{} % Change to "all" if they're all listed
						% \begin{compactdesc}
							% \item[]
						% \end{compactdesc}
				% \end{compactdesc}
			
				% \begin{compactdesc}
				
					% \item[BDAG ] \phantom{.}
						% \begin{compactitem}
							% \item 
						% \end{compactitem}
						
					% \item[Brill , PDF ] \phantom{.}
						% \begin{compactitem}
							% \item 
						% \end{compactitem}
						
					% \item[CGL , PDF ] \phantom{.}
						% \begin{compactitem}
							% \item 
						% \end{compactitem}
						
					% \item[LsJ , PDF ] \phantom{.}
						% \begin{compactitem}
							% \item 
						% \end{compactitem}
						
					% \item[TDNT ] \phantom{.}
						% \begin{compactitem}
							% \item 
						% \end{compactitem}
						
					% \item[TDNTA] \phantom{.}
						% \begin{compactdesc}
							% \item 
						% \end{compactdesc}
						
					% \item[NIDNTT] \phantom{.}
						% \begin{compactdesc}
							% \item 
						% \end{compactdesc}
						
					% \item[NIDNTTE] \phantom{.}
						% \begin{compactdesc}
							% \item 
						% \end{compactdesc}
						
				% \end{compactdesc}
		
	% -----
	\subsubsection{Book of Mormon} % *BOM
	% -----
		\label{MERIT-BOM}
		
		% --
		\paragraph{Noun}
		% --
	
			\begin{tabular}{ll}
				Total occurrences in the BOM & 5
			\end{tabular}
			
			\begin{compactdesc}
				\item[{\hyperref[2NEPHI-2-8]{2 Nephi 2:8}}]
				\item[{\hyperref[2NEPHI-31-19]{2 Nephi 31:19}}]
				\item[{\hyperref[ALMA-24-10]{Alma 24:10}}]
				\item[{\hyperref[HELAMAN-14-13]{Helaman 14:13}}]
				\item[{\hyperref[MORONI-6-4]{Moroni 6:4}}]
			\end{compactdesc}
			
		% --
		\paragraph{Verb}
		% --
	
			\begin{tabular}{ll}
				Total occurrences in the BOM & 2
			\end{tabular}
			
			\begin{compactdesc}
				\item[{\hyperref[MOsIAH-2-19]{Mosiah 2:19}}]
				\item[{\hyperref[ALMA-22-14]{Alma 22:14}}]
			\end{compactdesc}
		
	% -----
	\subsubsection{Doctrine and Covenants} % *DC
	% -----
		\label{MERIT-DC}
		
		% --
		\paragraph{Noun}
		% --
	
			\begin{tabular}{ll}
				Total occurrences in the DC & 1
			\end{tabular}
			
			\begin{compactdesc}
				\item[{\hyperref[DC3-20]{DC 3:20}}]
			\end{compactdesc}
			
		% --
		\paragraph{Verb}
		% --
		
			\begin{tabular}{ll}
				Total occurrences in the DC & 0
			\end{tabular}
		
	% % -----
	\subsubsection{Pearl of Great Price} % *PGP
	% -----
		\label{MERIT-PGP}
	
		\begin{tabular}{ll}
			Total occurrences in the PGP & 0
		\end{tabular}
		
		% \begin{compactdesc}
			% \item[]
		% \end{compactdesc}
		
% ----------------------
\subsection{Key Phrases} % *KEYPHRAsEs *PHRAsEs
% ----------------------
	\label{MERIT-PHRAsEs}

	\begin{compactdesc}
	
		\item[merit(s) of Christ] \textbf{6} | \hyperref[2NEPHI-2-8]{2 Nephi 2:8}; \hyperref[2NEPHI-31-19]{2 Nephi 31:19}; \hyperref[ALMA-24-10]{Alma 24:10}; \hyperref[HELAMAN-14-13]{Helaman 14:13}; \hyperref[MORONI-6-4]{Moroni 6:4}; \hyperref[DC3-20]{DC 3:20}
			% \begin{compactdesc}
				% \item[Bible anchor verses] \phantom{.}
					% \begin{compactdesc}
						% \item[Romans 5:5] \gr{}
					% \end{compactdesc}
				% \item[Commentary] ---
			% \end{compactdesc}
			
		% \item[] \textbf{\#times} | list
			% % \begin{compactdesc}
				% % \item[Bible anchor verses] \phantom{.}
					% % \begin{compactdesc}
						% % \item[Romans 5:5] \gr{}
					% % \end{compactdesc}
				% % \item[Commentary] ---
			% % \end{compactdesc}
			
	\end{compactdesc}
	
% % ----------------------
% \subsection{Bible Dictionaries} % *DICTIONARIEs
% % ----------------------
	% \label{MERIT-BD}

% ----------------------
\subsection{Restoration Usage} % *REsTORATION
% ----------------------
	\label{MERIT-REsTORATION}

	% -----
	\subsubsection{Joseph smith} % *Js
	% -----
		\label{MERIT-Js}
	
		\begin{compactdesc} % Format is: 21 May 1843 HC-a, or whatever date, whatever record-keeper (if there are multiple), then the passage 
			\item[{\hyperlink{EMs-1834-Js-CIR-MAR-1834-q}{Circa March 1834}}] But notwithstanding this trangression, by which man had cut himself off from an immediate intercourse with his Maker without a Mediator, it appears that the great and glorious plan of his redemption was previously meditated; the sacrifice prepared; the atonement wrought out in the mind and purpose of God, even in the person of the son, through whom man was now to look for acceptance, and through whose \tr{merits} he was now taught that he alone could find redemption, since the word had been pronounced, Unto dust thou shalt return!
		\end{compactdesc}
		
	% % -----
	% \subsubsection{Temple Ordinances}
	% % -----
		% \label{MERIT-TEMPLE}
	
		% \begin{compactdesc}
			% \item[{\hyperlink{KEY}{Date/source}}]
		% \end{compactdesc}
	
	% % -----
	% \subsubsection{Brigham Young} % *BY
	% % -----
		% \label{MERIT-BY}
	
		% \begin{compactdesc} % Format is: 21 May 1843 HC-a, or whatever date, whatever record-keeper (if there are multiple), then the passage
			% \item[{\hyperlink{KEY}{Date/source}}]
		% \end{compactdesc}
	
% ----------------------
\subsection{Commentary and studies} % *COMMENTARY *sTUDIEs
% ----------------------
	\label{MERIT-COMMENTARY}

	% -----
	\subsubsection{Key Points}
	% -----
		\label{MERIT-KEYs}
	
		\begin{compactitem}
			\item All six uses of the noun ``merit'' in the scriptures are plural; it's always the \textit{merits} of Christ, and not just the \textit{merit}.
			\item All six uses of the noun ``merit'' refer to Christ's \textit{merits}. None of the refer to the merits of people.
				\begin{compactitem}
					\item The two uses of the noun ``merit'' both involve people, not Christ---one is King Benjamin, and the other is humankind in general---and both are explicit that these people either do not merit anything of themselves, or that God merits much more.
				\end{compactitem} 
		\end{compactitem}
		
	% -----
	\subsubsection{Use of the word ``merit'' in the \textit{Book of Common Prayer}}
	% -----
		\label{MERIT-common}
		
		The King James version of the Bible does not have any uses of the term ``merit,'' either noun or verb. But the \textit{Book of Common Prayer}, a text in use in the Church of England in the 1600's, does use the term, and the restoration uses are very similar to how that book uses it.
		
		Here are the uses from the \textit{Book of Common Prayer} (page numbers are to PDF page numbers in the version I have):
		
			\begin{compactitem}
				\item ``through the merits of Jesus Christ our saviour'' (22)
				\item ``who art the only giver of all victory, through the merits of thy son Jesus Christ our Lord'' (32)
				\item ``we may pass to our joyful resurrection, for his merits who died, and was buried, and rose again for us'' (82)
				\item ``that we may alway serve thee in pureness of living and truth, through the merits of the same thy son Jesus Christ our Lord'' (87)
				\item ``through the merits of Christ Jesus our saviour'' (94)
				\item ``through the merits of Christ Jesus our saviour'' (95)
				\item ``through the merits of Christ Jesus our saviour'' (96)
				\item ``but through the merits and mediation of Jesus Christ thy son, our Lord'' (110)
				\item ``through the merits of Jesus Christ our Lord'' (110)
				\item ``through the merits of thy son Jesus Christ our Lord'' (140)
				\item ``that by the merits and death of thy son Jesus Christ, and through faith in his blood, we and all thy whole Church may obtain remission of our sins, and all other benefits of his passion'' (157)
				\item ``not weighing our merits, but pardoning our offences, through Jesus Christ our Lord'' (157)
				\item ``by the merits of the most precious death and passion of thy dear son'' (157)
				\item ``through the merits of thy most dearly beloved son Jesus Christ our Lord'' (183)
				\item ``through the merits and mediation of Jesus Christ, thine only son, our Lord and saviour'' (184)
				\item ``which may in the end bring us to life everlasting, through the merits of Jesus Christ, thine only son our Lord'' (185)
				\item ``through the merits and mediation of Jesus Christ our Lord'' (185)
				\item ``And after the multitude of thy mercies look upon us; Through the merits and mediation of thy blessed son, Jesus Christ our Lord'' (195)
				\item ``Hear, Lord, and save us, for the infinite merits of our blessed saviour, thy son, our Lord Jesus Christ'' (310)
				\item ``From this unnatural Conspiracy, not our merit, but thy mercy; not our foresight, but thy providence delivered us'' (316)
				\item ``But pardon us for thy mercies sake, through the merits of thy son Jesus Christ our Lord'' (320)
				\item ``Through the merits and mediation of thy blessed son Jesus Christ our Lord'' (321)
				\item ``Grant this for the all-sufficient merits of thy son our saviour Jesus Christ' (322)
				\item ``From all these, O gracious and merciful Lord God, not our merit, but thy mercy'' (326)
				\item ``by the merits and mediation of Christ Jesus our saviour'' (329)
				\item ``We are accounted righteous before God, only for the merit of our Lord and saviour Jesus Christ, by faith, and not of our own works or deservings'' (334)
			\end{compactitem} 